% Method sections for the revised paper: codebook ceiling + self-consistent
% quantization. Included from main.tex.

\section{What a Paired 8-Bit Weight Code Can Actually Deliver}
\label{sec:ceiling}

Our original codec pairs two real weights as $z=w_a+iw_b$ and spends three
magnitude bits and five phase bits per pair. Before asking whether that code
helps a language model, we ask a prior question the earlier drafts never posed:
at eight index bits per pair, how much room is there at all, and is the polar
arrangement a good way to use it?

\paragraph{A matched-strength control.} Every comparison in the previous draft
placed the polar code against a real 4-bit control whose codebook was a
\emph{uniform} 16-level grid on $[-1,1]$. The polar code, by contrast, gets an
implicitly non-uniform effective codebook, because a product of a radius grid
and a phase grid is not a square lattice. That asymmetry, not the phase
representation, can account for a small win. We therefore add a control that
gives the scalar code the same per-tensor shape adaptation: a 16-level
Lloyd--Max codebook fitted to the tensor's own normalized weight
distribution, with the identical per-row FP32 scale. The extra storage is 16
FP32 per tensor, which is under $10^{-4}$ bits per weight for these matrices.

\paragraph{An upper bound.} We also fit an \emph{unconstrained} 256-point
two-dimensional codebook to the pairs by $k$-means. Any scheme that pairs two
real weights and spends eight index bits on the pair -- polar, Cartesian,
rotated, or otherwise -- is a constrained special case of this code, so its
distortion lower-bounds the whole family at this budget.

\begin{table}[ht]
\centering\small
\setlength{\tabcolsep}{4pt}
\caption{Weight reconstruction at eight index bits per pair, relative to the Lloyd-fitted
real-4 control (values below $1.000$ beat it). \texttt{vq2d} is an unconstrained 256-point
2-D codebook and lower-bounds distortion for every paired code at this budget.
Every entry uses the best of four pairings and a per-row FP32 scale.}
\label{tab:ceiling}
\begin{tabular}{llrrrrr}\toprule
Checkpoint & Tensor & real4-Lloyd MSE & real4-unif. & polar-legacy & polar-fitted & vq2d\\\midrule
Qwen3.5-0.8B & DeltaNet QKV L0 & 3.252e-06 & 1.308 & 1.194 & 1.065 & \textbf{0.759}\\
Qwen3.5-0.8B & DeltaNet QKV L8 & 3.135e-06 & 1.214 & 1.156 & 1.064 & \textbf{0.793}\\
Qwen3.5-0.8B & DeltaNet QKV L16 & 5.039e-06 & 1.240 & 1.165 & 1.059 & \textbf{0.795}\\
Qwen3.8-27B & DeltaNet QKV L0 & 2.563e-06 & 1.186 & 1.162 & 1.073 & \textbf{0.818}\\
Qwen3.8-27B & Attention Q L3 & 3.317e-06 & 1.217 & 1.154 & 1.060 & \textbf{0.791}\\
Qwen3.8-27B & MLP up L3 & 1.182e-06 & 1.169 & 1.163 & 1.081 & \textbf{0.817}\\
\bottomrule
\end{tabular}
\end{table}


Three facts follow from Table~\ref{tab:ceiling}, and they hold on every tensor
we measured, across two checkpoints.

\begin{enumerate}
\item The reported polar advantage is an artifact of the control. Polar beats
the uniform scalar code, but the Lloyd-fitted scalar code beats polar on all
six tensors.
\item The polar family is structurally capped. Even after we repair its
assignment rule, fit its radius grid, and sweep every magnitude/phase bit split
(Section~\ref{sec:polarfix}), the best polar code remains $6$--$8\%$ above the
fitted scalar code.
\item The \emph{pairing} idea is nonetheless sound. A free 2-D codebook sits
$18$--$24\%$ \emph{below} the fitted scalar code at the same index budget. The
headroom is real; the polar parameterization simply fails to reach it.
\end{enumerate}

We therefore retain joint 2-D quantization of weight pairs and discard the
magnitude/phase parameterization. This also removes the practical obstacle that
defeated our packed kernels: decoding a polar index requires a sine and a
cosine per pair, which made decode-plus-GEMM $3$--$4.5\times$ slower than
resident BF16, whereas decoding a 2-D codebook index is a 256-entry table
lookup that fits in shared memory.

\subsection{Why repairing the polar assignment rule changes almost nothing}
\label{sec:polarfix}

The original quantizer rounds magnitude and phase independently. That is not
nearest-neighbour assignment in the $8\times32$ codebook. For a decoded phase
ray $\theta_p$ and a source $z=re^{i\theta}$,
\begin{equation}
 |z-g_m e^{i\theta_p}|^2 = r^2\sin^2\delta + (g_m - r\cos\delta)^2,
 \qquad \delta=\theta-\theta_p ,
\end{equation}
so the optimal radius level is the one nearest the \emph{projection}
$r\cos\delta$, not the one nearest $r$. We implemented exact nearest-neighbour
assignment over all 256 codepoints. It reduced weight MSE by $0.035\%$.

The reason is quantitative and worth stating, because it retrospectively
explains a long series of null results in this project. With five phase bits
the nearest-ray error satisfies $|\delta|\le\pi/32$, hence
$1-\cos\delta \le 4.8\times10^{-3}$. The projection correction is therefore at
most half a percent of the radius, while one step of the three-bit radius grid
is about $14\%$ of the row scale. The correction almost never changes which
radius level wins. Every phase-side refinement we tried across this project --
learned global phase offsets, activation-weighted phase choice, output-error
phase selection -- was operating on a quantity bounded to be small. The binding
constraint is the radius grid, not the phase grid, and the sweep over bit
splits confirms this: $8\times32$ is optimal, and moving bits toward the phase
($4\times64$, $2\times128$) is catastrophic.

\section{Self-Consistent Weight Quantization}
\label{sec:scq}

\subsection{The calibration fixed point}

Activation-aware post-training quantization chooses a code for layer $\ell$ by
minimizing a metric weighted by that layer's input statistics
$H_\ell=\mathbb E[x_\ell x_\ell^\top]$ (or its diagonal), estimated on
calibration data. GPTQ, AWQ, and every activation-weighted variant in this
project measure those statistics on the \emph{full-precision} model.

Once many layers are actually replaced, this is no longer the right statistic.
The inputs to layer $\ell$ are produced by layers $1,\dots,\ell-1$, which are
themselves quantized, so the deployed model presents layer $\ell$ with a
different distribution from the one its codebook was fitted to. Writing
$\mathcal Q$ for the per-layer quantizer and $H[\cdot]$ for the calibration
statistics a given set of weights induces, the conventional procedure computes
\begin{equation}
 Q_{\text{one-shot}} = \mathcal Q\bigl(H[W_{\text{BF16}}]\bigr),
\end{equation}
whereas the quantity that describes the deployed model is the fixed point
\begin{equation}
 Q^\star = \mathcal Q\bigl(H[Q^\star]\bigr).
 \label{eq:fixedpoint}
\end{equation}
The gap between them grows with the number of simultaneously quantized layers.
This is exactly the setting of our main experiment, which replaces eighteen
projections at once, and it is a plausible source of the instability that
defeated our earlier selectors: a rule tuned against one snapshot of
activations need not remain optimal once the snapshot moves.

\subsection{An SCF loop}

Equation~\eqref{eq:fixedpoint} is structurally the self-consistency condition of
Kohn--Sham density functional theory, and it admits the same iterative
treatment. Table~\ref{tab:dft} gives the correspondence.

\begin{table}[ht]
\centering\small
\caption{Correspondence between the Kohn--Sham SCF loop and self-consistent
quantization. The analogy is a numerical-methods one; no physical or quantum
claim is intended.}
\label{tab:dft}
\begin{tabular}{ll}\toprule
Kohn--Sham DFT & Self-consistent quantization\\\midrule
Electron density $n(r)$ & Activation statistics $H_\ell$\\
Effective potential $v_{\rm eff}[n]$ & Weighted per-layer quantizer metric\\
Solve KS equations $\to n'$ & Re-quantize, re-measure $\to H'$\\
Total energy $E[n]$ & Calibration NLL of the quantized network\\
Density mixing (linear, Pulay/DIIS) & Statistics mixing (linear, Anderson)\\
SCF convergence & Codebook-assignment churn $\to0$\\
Charge sloshing & Assignment oscillation between codepoints\\\bottomrule
\end{tabular}
\end{table}

One iteration is: measure $H^{(n)}_{\rm out}=H[Q^{(n)}]$ with a single forward
pass over the calibration windows; form the residual
$R^{(n)}=H^{(n)}_{\rm out}-H^{(n)}$; mix
\begin{equation}
 H^{(n+1)} = H^{(n)} + \alpha R^{(n)}
\end{equation}
or accelerate with Anderson/Pulay mixing over the last $m$ residuals; and
re-quantize every layer under $H^{(n+1)}$, warm-starting each codebook from the
previous iteration. Iteration zero is precisely the conventional one-shot
calibration, so the loop is a strict extension of current practice.

\subsection{Energy guard}

The analogy is not exact in one important respect. Kohn--Sham SCF iterates a
map whose fixed points are stationary points of a variational energy
functional; our assignment step is discrete and carries no such guarantee, so a
naive iteration can oscillate or converge to a worse code. We therefore treat
the calibration NLL as the total energy and guard every step: a trial iterate
is accepted only if it does not increase that energy, and otherwise it is
rejected and $\alpha$ is halved. Any Anderson history is discarded on a
rejection, as in a DFT code that restarts DIIS after a bad cycle.

This guard has a useful consequence. Because iteration zero \emph{is} the
one-shot solution and no accepted step increases the calibration energy, the
loop cannot terminate worse than conventional calibration on the calibration
set. The method is therefore a safe drop-in on top of any codec, which is how
we evaluate it: we apply the identical loop to the scalar, polar, and 2-D
codecs and report the improvement each one gains.

\paragraph{Cost.} An iteration costs one forward pass over the calibration
windows plus one re-quantization, and no backward pass. Our runs converge or
stall within a handful of iterations, so the total calibration cost stays in
the minutes for a 0.8B model. The deployed model is unchanged in size, format,
and decode cost: self-consistency moves only \emph{which} codepoints the
indices name, never how many bits they take.
